\documentclass[10pt,a4paper]{amsart}
\usepackage[margin=19mm]{geometry}
\usepackage{amsmath,amssymb,mathtools}
\usepackage[scr=rsfs,cal=euler]{mathalfa}
\usepackage{xfrac}
\usepackage[unicode,hidelinks,bookmarks=false]{hyperref}
\newcommand{\ie}{i.e.\ }
\newcommand{\cf}{cf.\ }
\newcommand{\resp}{respectively\ }
\newcommand{\Subsection}{\subsection}
\providecommand{\nobreakdash}{\nobreak\hbox{-}}
\setlength{\emergencystretch}{2em}
\multlinegap0pt
\usepackage{bm}
\usepackage{bbm}
\usepackage{dsfont}
\usepackage{xspace}
\usepackage{cancel}
\usepackage{mathtools}
\usepackage{amsthm}
\makeatletter
\@ifundefined{theorem}{\newtheorem{theorem}{Theorem}}{}
\@ifundefined{corollary}{\newtheorem{corollary}[theorem]{Corollary}}{}
\makeatother
\title[Ricci Solitons on a family of three dimensional Lorentzian W]{Ricci Solitons on a family of three dimensional Lorentzian Walker manifolds}
\author{A. Diatta, M. Ciss, A. S. Diallo}
\date{Source: arXiv:2508.10217v1 (CC BY 4.0)}
\begin{document}
\maketitle

\noindent\emph{Source and licence.} Ricci Solitons on a family of three dimensional Lorentzian Walker manifolds. Version arXiv:2508.10217v1 (CC BY 4.0), distributed under the Creative Commons Attribution 4.0 International licence, as recorded in the arXiv licence metadata for this exact version and shown on the abstract page https://arxiv.org/abs/2508.10217v1. Full rights notice: licenses/arxiv-2508.10217-CC-BY-4.0.txt.\par\smallskip

\noindent\textit{Editorial scope.} Counted unit: the Corollary whose source label is \texttt{eq3.8} in the article (source0.tex lines 371--392, source label \texttt{eq3.8}), delivered together with the article's own complete original proof of it (source0.tex lines 394--462, one \texttt{proof} environment of 69 lines). No mathematics is written, repaired or re-derived: statement and proof are the article's own text line for line. Required context retained uncounted: eq2.1 (lines 122--129); eq3.1 (lines 179--181); eq3.3 (lines 215--242). Every definition, display and label of those lines is retained unchanged. Omitted from this package and not redistributed: 1--121, 130--178, 182--214, 243--370, 393--393, 463--601. Nothing inside a retained excerpt is omitted or altered: every retained body is delivered exactly as the article's lines are written, so no omission receipt of any kind accompanies this package. Citations: the retained excerpts carry no citation command at all, so no bibliography entry of the article is transcribed and no reference is redistributed. Reference adaptation: none was needed and none was made; every reference inside the delivered unit resolves inside it, so no unresolved reference or citation is produced. Printed slips kept literal: no printed word, symbol, hypothesis or number of the counted statement or of its proof was altered. Layout and class: the article's own class is \texttt{amsart} with packages including amsmath, amssymb, amsthm, times, latexsym, eucal; the wrapper is the lane's pinned \texttt{amsart} editorial wrapper at 10pt on A4, which restates the theorem-like environments the retained text needs and adds the article's own macro definitions that the retained text needs (none), with extra wrapper packages (bm, bbm, dsfont, xspace, cancel, mathtools). The delivered numbering is the unit's own. Assets: no retained span contains a figure environment, a graphics-inclusion command, a PDF-page inclusion, a table or a TikZ picture; no image file is copied into this package, the package declares no asset dependency and no image file is redistributed with it. Licence: version arXiv:2508.10217v1 of this article is distributed under the Creative Commons Attribution 4.0 International licence, as recorded in the arXiv licence metadata for this exact version and shown on the abstract page https://arxiv.org/abs/2508.10217v1. A full rights notice is delivered at licenses/arxiv-2508.10217-CC-BY-4.0.txt. Characters: every retained excerpt is pure ASCII, so the delivered engine, XeLaTeX with its default Latin Modern text font, reports no missing character.
\begin{eqnarray}\label{eq2.1} 
g _f= \left( \begin{array}{cccc}
0&0&1\\
0&\varepsilon&0\\
1&0&f(t,x,y)
\end{array}
\right),
\end{eqnarray}

\begin{eqnarray}\label{eq3.1} 
\mathcal{L}_X g + \rho = \lambda \cdot g,
\end{eqnarray}

\begin{theorem}
Let $(M, g_f)$ be a three-dimensional Walker manifold, where $g_f$ is the metric 
given by \eqref{eq2.1} and $f$ is the defining function given by 
$f(t,x,y) = a(t,x) y^{2} + b(t,x) y + d(t,x)$.  Then $(M, g_f, X, \lambda)$ is a Ricci soliton 
if equation \eqref{eq3.1} is satisfied by $X = A\partial_t+ B\partial_x + C \partial_y $ 
where $A, B $ and $C$ are given by :
\begin{eqnarray}\label{eq3.3}
A(t,x, y) &=& \left( \lambda-K_y(y)\right)t + \varepsilon H_{y}(y)
- \frac{1}{2}a_{t}(t)y^2 +N(x,y), \nonumber\\ 
B(t,x,y) &=&\frac{\lambda}{2}x + H(t,y),\nonumber\\
C(t,x,y)&=& -\varepsilon H_{t}(y)x+K(y),
\end{eqnarray}
and the functions $H, K, E$ satisfy the following conditions :
\begin{align*}
&\left( \frac{1}{2}a_{tx}(t,x)	-\varepsilon aH_{t}(t,y)\right) y^{2}
+ \left(  \frac{1}{2}b_{tx}(t,x)-\varepsilon bH_{t}(t,y)\right) y\nonumber\\
& + \frac{1}{2}d_{tx}(t,x)-\varepsilon dH_{t}(t,y)+\varepsilon H_{y}(t,y) + N_{x}(x,y) = 0,
\end{align*} 
and
\begin{align*}
&&\left(-\varepsilon H_{t}(y)x + K(y)\right) f_y + f\left( -2\varepsilon H_{ty}(y)x
+ 2K_y(y) + \lambda + \frac{1}{2}f_{tt}\right)\nonumber  \\
&& + \left( ( \lambda - tK_y(y)) + \varepsilon H_{y}(y) - \frac{1}{2}a_{t}(t)y^2 
+ N(x,y)\right) f_t\nonumber\\
&& + 4\left(\varepsilon H_{yy}(y) + N_y(x,y) - tK_{yy}(y) - a_ty \right) 
- \frac{1}{2\varepsilon}f_{xx}=0.
\end{align*}
\end{theorem}

 \begin{corollary}
Let $(M, g_f)$ be a flat Walker manifold three-dimensional, where $g_f$ 
denotes the metric defined by (\ref{eq2.1}) and $f$ is the defining function 
given by $f(y)=\alpha y^{2}+\beta y+\gamma$. Then$(M, g_f, X, \lambda)$ 
is a Ricci soliton  if equation \eqref{eq3.1} is satisfied by 
$X = A\partial_t+ B\partial_x + C \partial_y $ where $A, B $ and $C$ are given by :
\begin{eqnarray}\label{eq3.8}
A(t,x, y) &=& - \frac{\beta}{2} \left( \frac{\lambda}{4}y+ \delta\right)y  
		+ \frac{\lambda }{2} t +F(x,y), \nonumber\\ 
B(t,x,y) &=&\frac{\lambda}{2}x+H(t,y),\nonumber  \\ 
C(t,x,y) &=& -\varepsilon H_{t}(y)x+K(y),
\end{eqnarray}
and the functions $H,K,F$ satisfy the following conditions :
\begin{equation*}
2H_y(y) - \varepsilon fH_t(y) + F_x(x,y)=0,
\end{equation*}
and
\begin{eqnarray*}
\left(K(y) - \varepsilon H_{t}(y)x\right)f_y &+& 2\varepsilon H_{yy}(y)x - 2K_{yy}(y)t \\
	&+& 2F_{y}(x,y) + f\left( 2K_{y}(y) + \lambda\right) =0. 
\end{eqnarray*}
\end{corollary}

\begin{proof}
Let $f(y)=\alpha y^{2}+\beta y+\gamma$ and 
$X = A\partial_t+ B\partial_x + C \partial_y $ in $(M, g_f)$.  Using (\ref{eq2.1}),  \eqref{eq3.1} and (\ref{eq3.3}), a standard calculation gives that three-dimensional 
Walker manifold $(M,g_f)$ is a Ricci soliton if and only if the following conditions hold :
\begin{itemize}
\item (C''1) $C_t= 0$,  
\item (C''2) $\varepsilon B_t +C_x= 0$,  
\item (C''3) $A_t +C_y +fC_t=\lambda$,  
\item (C''4) $2\varepsilon B_x=\varepsilon\lambda$,  
\item (C''5) $\varepsilon B_y +fC_x +A_x=0$,  
\item (C''6) $ X(f)+2\left( fC_y+2A_y\right)  =\lambda f$. 
\end{itemize}
From condition (C''1), we have: $C(t,x,y) = C(x,y)$. We establish the form of $B$ 
under condition (C"4) :
\begin{equation}\label{eq3.9}
B(t,x,y)= \frac{\lambda}{2}x + H(t,y),
\end{equation}
with $H \in C^{\infty}(M)$. We compute $B_t$ and substitute it into condition (C"2), 
we obtain : $\varepsilon H_{t}(t,y) + C_{x}=0$. Therefore
\begin{equation}\label{eq3.10}
	C(t,x,y)=-\varepsilon H_{t}(y)x + K(y),
\end{equation}
where $K\in C^{\infty}$. Using \eqref{eq3.10} in condition (C"3), we have :
\begin{eqnarray*}
A_t -\varepsilon H_{ty}(y)x + K_{y}(y) = \lambda,
\end{eqnarray*} 
and then
\begin{equation}\label{eq3.11}
A(t,x,y)=\varepsilon H_{y}(y)x+\left( \lambda-K_{y}(y)\right) t + F(x,y),
\end{equation}
where $F\in C^{\infty}(M)$. \\

We differentiate conditions (C"3) and (C"5) respectively with respect to $t$ and $x$, 
we get :
\begin{eqnarray*}
A_{tx} + C_{yx}=0 \quad \mbox{and}\quad
A_{xt} + f_{t}C_{x}+\varepsilon B_{yt} = 0.
\end{eqnarray*}
Hence
\begin{eqnarray*}
- C_{yx} + f_tC_x + \varepsilon B_{yt} = 0.
\end{eqnarray*}
Using \eqref{eq3.9} and \eqref{eq3.10}, we have : $ H_{yt}(y)=0$ and
so $H_{y}(y,t) = H_{y}(y)$. Finally we have :
\begin{eqnarray*}
A(t,x,y) &=& \varepsilon H_{y}(y)x + \left( \lambda - K_{y}(y)\right) t + F(x,y),\\
B(t,x,y) &=& \frac{\lambda}{2}x + H(t,y),\\
C(t,x,y) &=& - \varepsilon H_{t}(y)x+K(y).
\end{eqnarray*}
Thus, condition $(C"5)$, gives us
\begin{equation*}
2H_y(y) - \varepsilon f H_t(y) + F_x(x,y)=0.
\end{equation*}
Using the condition (C"6), we obtain :
\begin{eqnarray*}
X(f) + 2\varepsilon H_{yy}(y)x - 2K_{yy}(y)t + 2F_{y}(x,y) 
+ f\left( 2K_{y}(y)+\lambda\right) &=& 0.
\end{eqnarray*}
Compute $X(f)$, we have :
\begin{eqnarray*}
X(f) = \left[ -\varepsilon H_{t}(y)x + K(y)\right] f_y. 
\end{eqnarray*}
Then
\begin{eqnarray*}
\left( -\varepsilon H_{t}(y)x + K(y)\right)f_y &+& 2\varepsilon H_{yy}(y)x 
-2K_{yy}(y)t \\
&+& 2F_{y}(x,y) + f\left( 2K_{y}(y) + \lambda\right) =0.   
	\end{eqnarray*}
\end{proof}	

\end{document}
